\documentclass[12pt,a4paper]{article}
\usepackage{amsmath}
\usepackage{amssymb}
\usepackage{amsthm}
\usepackage{geometry}
\usepackage{xcolor}
\usepackage{hyperref}

\geometry{margin=1in}

% Define theorem environments
\theoremstyle{definition}
\newtheorem{definition}{Definition}[section]
\newtheorem{theorem}{Theorem}[section]
\newtheorem{proposition}{Proposition}[section]
\newtheorem{lemma}{Lemma}[section]
\newtheorem{remark}{Remark}[section]
\newtheorem{example}{Example}[section]

\title{COMP0085: Approximate Inference}
\author{Xu Chen}
\date{\today}

\begin{document}

\maketitle
\tableofcontents
\newpage

\section*{Introduction}

These lecture notes cover approximate inference methods for probabilistic graphical models. The course develops techniques for computing with complex distributions where exact inference is intractable.

\vspace{0.5cm}
\noindent\textbf{Course Structure:}
\begin{enumerate}
    \item \textbf{Foundation (Sections 1-3):} Graphical models, conditional independence, exact inference algorithms (polytrees, junction trees)
    \item \textbf{Approximate Inference Methods (Sections 4-9):} Variational inference, expectation propagation, loopy belief propagation
    \item \textbf{Theoretical Foundations (Sections 10-11):} Bethe free energy, convex duality, connection to statistical physics
    \item \textbf{Applications and Extensions (Sections 12-15):} Model selection, Gaussian processes, Boltzmann machines, deep learning approaches
\end{enumerate}

The core mathematical framework builds progressively: from exact inference on trees, to approximate inference on general graphs, to the theoretical justification via variational bounds and convex optimization.

\newpage
\vspace*{2cm}
\begin{center}
    \section*{\Large Part I: Foundations --- Graphical Models and Exact Inference}
\end{center}
\vspace*{1cm}
\nobreak

The foundation of this course rests on understanding graphical models and their induced conditional independence structures. We begin with exact inference algorithms on trees, which serve as the basis for understanding approximate inference on graphs with loops.

\newpage

\section{Graphical Models and Conditional Independence}

\subsection{Factor Models}

\begin{definition}[Factor Model]
A factor model represents a probability distribution as a product of factors:
\[
P(Z) = \frac{1}{Z} \prod_{i=1}^m f_i(Z_i)
\]

where factors are functions of subsets of variables. Each factor $f_i$ induces conditional dependencies between variables in its scope.
\end{definition}

\subsection{Conditional Independence}

\begin{definition}[Conditional Independence in Graphs]
Variables $X$ and $Y$ are conditionally independent given $V$ (written $X \perp Y | V$) if:
\[
P(X, Y | V) = P(X | V) P(Y | V)
\]
Equivalently, $P(X | Y, V) = P(X | V)$. In graphical models, this holds when every path between $X$ and $Y$ contains some variable in $V$ (for undirected graphs or specific graph types).
\end{definition}

\subsection{Neighbourhood Definition}

\begin{remark}[Neighbourhood Relationship in Graphical Models]
In factor models, variables are neighbours if they share a common factor, meaning there exists a factor that references both variables.
\end{remark}

\subsection{Markov Property}

\begin{definition}[Markov Property]
Variable $X$ is conditionally independent of all non-neighbours given its neighbours:
\[
X \perp Y \mid \text{ne}(X), \quad \forall Y \notin \{\text{neighbours of } X\}
\]

This is conditional independence of all non-neighbours given its neighbours.
\end{definition}

\subsection{Markov Blanket}

\begin{definition}[Markov Blanket]
The Markov blanket (or Markov boundary) of variable $X$ is a minimal set of variables $M$ such that $X$ is conditionally independent of all other variables given $M$. That is:
\[
X \perp \{\text{all other variables}\} | M
\]
In graphical models, the Markov blanket typically includes the variable's parents, children, and co-parents (parents of children).
\end{definition}

\subsection{Undirected Graphical Models}

\begin{definition}[Undirected Graphical Models]
Markov networks (or undirected graphical models) represent factors as cliques on the graph. Neighbours in a Markov net share a factor.

\textbf{Joint probability factors over the maximal cliques} (or complete subgraphs) $C$ of the graph.
\end{definition}

\subsection{Cliques and Maximal Cliques}

\begin{definition}[Cliques and Maximal Cliques]
Cliques are fully connected subgraphs (every two nodes are connected). Maximal cliques are cliques not contained in other cliques (cannot be enlarged to add more vertices).
\end{definition}

\subsection{Directed Acyclic Graphical (DAG) Model}

\begin{definition}[Directed Acyclic Graph (DAG) Model]
The joint probability factorizes:
\[
P(X_1, \ldots, X_n) = \prod_{i=1}^n P(X_i | X_{\text{pa}(i)}) \quad \text{pa}(i) \text{ are parents of node } i
\]

These are Bayesian networks.
\end{definition}

\subsection{Bayes-Ball Algorithm}

\begin{remark}[Bayes-Ball Algorithm]
The Bayes-Ball algorithm determines conditional independence in DAGs by asking: can a ball travel from $X$ to $Y$ without being blocked by conditioning set $V$? This determines $X \perp Y | V$.

\textbf{Rules:}
\begin{itemize}
    \item \textbf{Nodes not in } $V$: Balls pass through (either direction on chain)
    
    \item \textbf{Nodes in } $V$: Block balls traveling along edges involving two children or arrows into the node
    
    \item \textbf{At v-structures} ($X \leftarrow Z \rightarrow Y$): Block balls unless $Z$ or its descendants are in $V$
\end{itemize}

\textbf{Blocking set:} A path from $X$ to $Y$ is blocked by $V$ if:
\begin{itemize}
    \item There exists a non-collider node $C$ on the path with $C \in V$, or
    \item There exists a collider node $C$ on the path with $C \notin V$ and no descendant of $C$ is in $V$
\end{itemize}

If no unblocked path exists from $X$ to $Y$ given $V$, then $X \perp Y | V$.
\end{remark}

\subsection{Markov Boundary for X in DAG}

\begin{remark}[Markov Boundary in Directed Acyclic Graphs]
In a DAG, the Markov boundary of $X$ consists of: parents, children, and co-parents of children.
\[
\text{MB}(X) = \text{pa}(X) \cup \text{ch}(X) \cup \text{pa}(\text{ch}(X))
\]

This is the minimal set such that $X \perp \text{(all other variables)} | \text{MB}(X)$.
\end{remark}

\subsection{Equivalence between Graphs and Families of Distributions}

\begin{remark}[Graphs Induce Families of Distributions]
A graph $G$ induces a family of distributions by defining conditional independence statements (CIS). For any graph:
\begin{itemize}
    \item \textbf{Directed graphs (DAGs):} $P(Z) = \prod_i P(X_i | \text{pa}(i))$ — global Markov implies factorization
    \item \textbf{Undirected graphs:} By Hammersley-Clifford theorem, positive distributions satisfying all global Markov properties factor as: $P(Z) = \frac{1}{Z} \prod_c f_c(Z_{C_c})$ where cliques $C_c$ correspond to maximal complete subgraphs
    \item \textbf{Factor graphs:} Can represent families not capturable by undirected graphs alone (more expressive)
\end{itemize}

All conditional independence statements implied by a graph $G$ are preserved in the corresponding family of distributions.
\end{remark}

\subsection{Adding Edges Changes Expressivity}

\begin{remark}[Effect of Adding Edges to Graphs]
Adding edges to a graph $G$ removes conditional independence assertions from the induced family:
\[
G' = G + \text{edge} \Rightarrow \mathcal{P}_{G'} \supseteq \mathcal{P}_{G}
\]

The new family is larger since we've removed constraints. This means denser graphs are more expressive but less constrained by conditional independence.
\end{remark}

\subsection{Polytrees}

\begin{definition}[Polytrees as Tree-Structured DAGs]
Polytrees are tree-structured DAGs that may have more than one root.
\[
P(Z) = \prod_i P(X_i | X_{\text{pa}(i)}) = \prod_i f_i(X_i)  \quad C_i = \cup \text{pa}(i) \quad \text{and} \quad f_i(X_{C_i}) = P(X_i | X_{\text{pa}(i)})
\]
\end{definition}

\subsection{Undirected Trees}

\begin{definition}[Undirected Trees and Pairwise Factors]
For undirected tree, all maximal cliques are of size 2; factor graph has pairwise factors.
\[
P(Z) = \frac{1}{Z} \prod_{e \in E} f_{uv}(X_i, Y_j)
\]
\end{definition}

\subsection{Single Rooted Directed Trees}

\begin{definition}[Single Rooted Directed Trees]
Single rooted directed tree can be written as a product of pairwise factors $\Rightarrow$ undirected tree.
\[
P(Z) = P(X_r) \prod_{i \neq r} P(X_i | X_{\text{pa}(i)}) = \prod_{e \in E} f_{uv}(X_i, X_r)
\]
\end{definition}

\newpage

\section{Polytrees, Message Passing, and Factor Models}

\subsection{Polytrees to Factor Models}

\begin{remark}[Polytrees and Factor Models]
A polytree (CDAG - directed acyclic graph) can be factorized as a factor model. A single rooted tree (CDAG) factorizes through pairwise factors via an undirected tree representation:
\[
P(X) = P(X_i) \prod_{j \in \text{ne}(i)} P(X_j | X_i) \text{ from belief propagation}
\]

An undirected tree can be obtained from the CDAG by treating edges as undirected.
\end{remark}

\subsection{Converting Undirected Trees to Rooted Directed Trees}

\begin{remark}[Root Selection and Directed Conversion]
Choose an arbitrary node $X_r$ to be the root and orient all arrows away from it. Then:

\begin{itemize}
    \item Compute $P(X_i)$ and $P(X_i | X_j)$ relationships
    \item The joint distribution becomes: 
    \[
    P(X_i) = P(X_r) \prod_{i \in \text{pa}(i)} P(X_i | X_{\text{pa}(i)}) = P(X_{r}) \prod_i \frac{P(X_i, X_j)}{P(X_j)} = \frac{\prod_i \text{numerator products}}{Z}
    \]
    where degree accounting: $\deg(i) = \text{\# being parent} + \text{\# being child}$
    \item For the rooted representation: $\deg(i) = \text{\# being parent}$, leaving one as prior
\end{itemize}
\end{remark}

\subsection{Finding Marginals in Undirected Trees}

\begin{remark}[Marginal Computation via Message Passing]
For an undirected tree with pairwise factorization:
\[
P(X) = \frac{1}{Z} \prod_{i,j \in E} f_{ij}(X_i, X_j)
\]

The marginal $P(X_i)$ is obtained by summing over all other variables. Using message passing, we compute:
\[
P(X_i) = \sum_{X_{\setminus i}} P(X) \propto \sum_{X_{\setminus i}} \prod_{j \in \text{ne}(i)} f_{ij}(X_i, X_j)
\]

This can be decomposed recursively. For each neighbour $j$ of $i$, define the message:
\[
M_{j \to i}(X_i) = \sum_{X_j} f_{ij}(X_i, X_j) \prod_{k \in \text{ne}(j) \setminus i} f_{jk}(X_j, X_k)
\]

The marginal factorizes as:
\[
P(X_i) \propto \prod_{j \in \text{ne}(i)} M_{j \to i}(X_i)
\]

Similarly, pairwise marginals are:
\[
P(X_i, X_j) = \sum_{X_{\setminus i,j}} P(X) \propto f_{ij}(X_i, X_j) M_{i \to j}(X_j) M_{j \to i}(X_i)
\]
\end{remark}

\subsection{Messages from Observed Leaf Nodes}

\begin{remark}[Conditioning on Observations]
When observing evidence, messages are conditioned rather than marginalized:

\begin{itemize}
    \item To compute $P(X_i)$ with no observations: 
    \[
    M_{a \to i}(X_i) = \sum_{X_a} f_{aX}(X_a, X_i) \propto \sum_{X_a} P(X_a | X_i) P(X_i)
    \]
    
    \item To compute $P(X_i | X_a = a)$ when $X_a$ is observed:
    \[
    M_{a \to i}(X_i) = f_{aX}(X_a=a, X_i) = \sum_{X_a} \delta_a(X_a) f_{aX}(X_a, X_i) \propto P(X_a = a | X_i) P(X_i)
    \]
\end{itemize}
\end{remark}

\subsection{Observed Internal Nodes}

\begin{remark}[Partitioning by Observed Nodes]
When an observed node is in the interior of the tree, it partitions the graph, and messages propagate independently:
\[
M_{b \to j} = f_{bX}(X_b = b, X) \quad \text{vs} \quad M_{a \to *} = f_{aX}(X_a = b, X)
\]
\end{remark}

\subsection{General Message Formulation}

\begin{remark}[Messages and Likelihood Scaling]
Messages $M_{i \to j}(X_j)$ are proportional to the likelihood of observed variables (O) within the subtree $T_{i \to j}$, possibly scaled by prior factors:
\[
M_{i \to j}(X_j) \propto P(X_{T_{i \to j}} \cap [O] | X_j) P(X_j)
\]

where the message encodes belief about $X_j$ given all evidence in the subtree rooted at $i$ (when viewing $j$ as the root).
\end{remark}

\subsection{Belief Propagation for Latent Class Models}

\begin{definition}[BP Updates for Latent Class Models]
For models with latent class variables, belief propagation follows:

\begin{itemize}
    \item Define $\alpha_i(k) = P(X_1, \ldots, X_i, \theta=k)$
    \item Forward pass: $\alpha_{i+1}(k) = P(X_{i+1} | \theta=k) \sum_k \alpha_i(k) \beta_i(k)$
    \item Backward pass: compute messages and marginals
\end{itemize}
\end{definition}

\subsection{Belief Propagation for Polytrees}

\begin{remark}[Factor Nodes in Polytrees]
Associated undirected tree contains loops created by non-binary latent classes. We can treat factors as message-passing nodes:
\[
P(X) = \prod_{i} \text{node}_i(X_i) \prod_{a} \text{factor}_a(X_a)
\]

Messages pass between variable nodes and factor nodes:

\begin{itemize}
    \item \textbf{Incoming} (node $\to$ factor): 
    \[
    M_{i \to a}(X_i) = \sum_{X_i} P(X_i | X_{\text{pa}(i)}) \prod_{b \neq a} M_{b \to i}(X_i)
    \]
    
    \item \textbf{Outgoing} (factor $\to$ node): 
    \[
    M_{a \to i}(X_i) = \frac{1}{Z_i} f_a(X_i, X_a) \prod_{j \in \text{ne}(a) \setminus i} M_{j \to a}(X_j)
    \]
\end{itemize}

This allows exact inference on polytree-structured graphs through efficient message passing.
\end{remark}

\newpage

\section{Junction Trees and Message Passing Algorithms}

\subsection{Converting Graphs While Preserving Distributions}

\begin{remark}[Equivalence Under Graph Transformation]
A transformation from graph $G$ to an easier-to-handle graph $G'$ is valid only if $P(X)$ can also be represented by $G'$. The critical principle:
\[
\text{Only remove conditional independence assertions; never add them.}
\]
\end{remark}

\subsection{Junction Tree Algorithm}

\begin{remark}[Construction Steps for Junction Trees]
The junction tree algorithm transforms a factor graph into a tree structure suitable for efficient inference:

\begin{enumerate}
    \item \textbf{DAGs to factor graphs}: Convert $P(X) = \prod P(X_i | X_{\text{pa}(i)}) = \prod f_i(X_{\alpha_i})$ where $\alpha_i = i \cup \text{pa}(i)$ and $P(X_r)$ is absorbed into an adjacent factor.
    
    \item \textbf{Incorporating observations}: For observed data $O$, the distribution becomes $P(X | O) \propto P(X, O) = P(X) \mathbf{1}_S(O)$. Add singleton factors with $\delta$-functions: $P(X, O) = \prod [f_a] f_b(O) f_c(C)$
    
    \item \textbf{Moralization}: Convert factor graph to undirected graph by replacing each factor with an undirected clique. Construct potentials on maximal cliques by multiplying factor potentials, ensuring each appears exactly once.
    
    \item \textbf{Triangulation}: Add edges so every cycle of length $\geq 4$ has at least one chord. Use either:
    \begin{itemize}
        \item \textit{Minimum deficiency search}: Choose variable inducing fewest new edges
        \item \textit{Maximum cardinality search}: Choose variable with most previously visited neighbours
    \end{itemize}
    The resulting chord graph must be direct (undirected).
    
    \item \textbf{Clique graph to junction tree}: Create nodes representing cliques; edges represent separators (intersections of cliques). This forms a tree structure.
\end{enumerate}
\end{remark}

\subsection{Running Intersection Property}

\begin{definition}[Running Intersection Property]
If two cliques contain variable $X$, all cliques and separators on the path between them must also contain $X$. This property is guaranteed if and only if the graph is chordal (triangulated).
\end{definition}

\subsection{Message Passing on Junction Trees}

\begin{remark}[Junction Tree Message Passing]
In a junction tree, nodes represent maximal cliques $C_i$, and edges represent separators. Messages pass between adjacent cliques, implementing distributional updates.

For a simple chain: $(ABC) - (BC) - (BCD)$ with separators $(B,C)$ and $(B,C,D)$, the message from $i$ to $j$ is:
\[
M_{i \to j}(X_{s_{ij}}) = \sum_{X_{\alpha_i} \setminus X_{s_{ij}}} f(C_i) \prod_{l \in \text{neigh}(i) \setminus j} M_{l \to i}(X_{s_{il}})
\]

The marginals are computed as:
\[
P(X_{\alpha_i}) \propto f_i(X_{\alpha_i}) \prod_{l \in \text{neigh}(i)} M_{l \to i}(X_{s_{il}})
\]

For separator marginals:
\[
P(X_{s_{ij}}) \propto M_{i \to j}(X_{s_{ij}}) M_{j \to i}(X_{s_{ij}})
\]
\end{remark}

\subsection{Local and Global Consistency}

\begin{remark}[Consistency Conditions]
If $g_i(X_{c_i})$ and $r_{ij}(X_{s_{ij}})$ are distributions satisfying the marginal consistency:
\[
\sum_{X_{\alpha_i} \setminus X_{s_{ij}}} g_i(X_{c_j}) = r_{ij}(X_{s_{ij}})
\]

then
\[
P(X) = \frac{\prod g_i(X_{c_i})}{\prod_{l \in \text{neigh}(i)} r_{ij}(X_{s_{ij}})}
\]
is also a valid distribution. The consistency is maintained because each separator potential must be counted only once. We have:
\[
g_i(X_{c_i}) = \sum P(X), \quad b_{ij}(X_{s_{ij}}) = \sum P(X)
\]
\end{remark}

\subsection{Hugin Propagation Algorithm}

\begin{remark}[Hugin Propagation Updates]
The Hugin algorithm maintains consistency through iterative updates:

\begin{itemize}
    \item \textbf{Initialization}: Set $\psi_i(X_{c_i}) \propto f_i(X_{c_i})$ and $r_{ij}(X_{s_{ij}}) \propto 1$ (unfactored); then $P(X) \propto \frac{\prod \psi_i(X_{\alpha_i})}{\prod r_{ij}(X_{s_{ij}})}$
    
    \item \textbf{Update for $i \to j$}: 
    \[
    \psi_i'(X_{s_{ij}}) = \sum q_i'(X_{c_i}), \quad r_{ij}'(X_{s_{ij}}) = \frac{\psi_i'(X_{s_{ij}}) r_{ij}(X_{s_{ij}})}{r_{ij}^{\text{old}}(X_{s_{ij}})}
    \]
    
    \item $P(X)$ remains unchanged after updates — local consistency is maintained throughout.
\end{itemize}
\end{remark}

\subsection{Computational Complexity}

\begin{remark}[Time and Memory Costs]
The computational cost of message passing is proportional to:
\[
O\left(\sum_i |X_{c_i}| \cdot |X_{c_i}|\right) = O\left(\sum_i \text{number of combinations of states in clique } i\right)
\]

This scales exponentially with the clique size, making triangulation heuristics important for finding small-treewidth decompositions.
\end{remark}

\subsection{Learning in Graphical Models}

\begin{remark}[Parameter Estimation via Expectation-Maximization]
For learning shared parameters $\theta$, the likelihood is:
\[
\ell = P(X_{\text{obs}} | \theta) = \sum_{X_{\text{latent}}} P(X_{\text{obs}}, X_{\text{latent}} | \theta)
\]

The log-likelihood can be decomposed:
\[
\log P(X_{\text{obs}}, X_{\text{latent}} | \theta) = \sum_j \log f_i(X_{c_j} | \theta) - \log Z(\theta) + H(\theta)
\]

where $Z(\theta)$ is the partition function. The EM algorithm iterates between:
\begin{itemize}
    \item \textbf{E-step}: Compute expected sufficient statistics under the posterior
    \item \textbf{M-step}: Update $\theta$ to maximize expected log-likelihood
\end{itemize}
\end{remark}

\newpage
\vspace*{2cm}
\begin{center}
    \section*{\Large Part IV: Applications and Extensions}
\end{center}
\vspace*{1cm}
\nobreak

The approximate inference methods developed in Part II have broad applicability. We now examine concrete applications: model comparison and selection, Gaussian processes as a nonparametric approach to regression, Boltzmann machines as classical models of distributed computation, and modern deep learning approaches that combine neural networks with variational inference.

\newpage

\section{Model Selection and Bayesian Inference}

\subsection{Bayesian Evidence}

\begin{definition}[Bayesian Model Evidence]
For comparing models $m$, the evidence is:
\[
P(D | m) = \int_{\theta_m} P(D | \theta_m, m) P(\theta_m | m) d\theta_m
\]

In conjugate prior settings:
\[
P(D | m) = \frac{\text{posterior value}}{\text{prior volume}}
\]
\end{definition}

\subsection{Laplace Approximation}

\begin{remark}[Laplace Approximation for Model Evidence]
To approximate $\log P(D, \theta_m | m)$, expand to second order around the maximum $\theta^*$ (the posterior mode):

\[
\log P(D, \theta_m | m) \approx \log P(D, \theta_m^* | m) + \nabla \log P(D, \theta_m^* | m)(\theta_m - \theta_m^*) + \frac{1}{2}(\theta_m - \theta_m^*)^T A (\theta_m - \theta_m^*)
\]

where $A := -\nabla^2 \log P(D, \theta_m^* | m)$ is the negative Hessian evaluated at the posterior mode.

At the mode, $\nabla \log P(D, \theta_m^* | m) = 0$, so:
\[
P(D | m) = \int e^{\log P(D, \theta_m | m)} d\theta_m \approx P(D | \theta_m^*, m) P(\theta_m^* | m) (2\pi)^{d/2} |A|^{-1/2}
\]

This approximates the posterior with a Gaussian.
\end{remark}

\subsection{Bayesian Information Criterion (BIC)}

\begin{remark}[Asymptotic Approximation with BIC]
The model evidence becomes:
\[
\log P(D | m) \approx \log P(\theta_m^* | m) + \log P(D | \theta_m^*, m) + \frac{d}{2} \log 2\pi - \frac{1}{2} \log |A|
\]

As $N \to \infty$, we have $A \to N A_0 + \text{const}$ where $A_0$ is a fixed positive definite matrix. Thus:
\[
\log |A| \to \log |N A_0| = d \log N + \log |A_0|
\]

Keeping only terms varying with $N$:
\[
\log P(D | m) \approx \log P(D | \theta_m^*, m) - \frac{d}{2} \log N
\]

This leads to model selection criteria:
\begin{itemize}
    \item \textbf{BIC}: $\text{BIC}' = d \log N - 2 \log P(D | \theta_m^*, m)$
    \item \textbf{AIC}: $\text{AIC} = 2d - 2 \log P(D | \theta_m^*, m)$
\end{itemize}

We can use ML estimate of $\theta$ instead of MAP for these approximations.
\end{remark}

\subsection{Automatic Relevance Determination (ARD)}

\begin{remark}[Hyperparameter Learning via ARD]
For prior $\alpha \propto N(0, C)$ with $\psi_i \sim N(W^T X_i, \sigma^2)$, and posterior on $W$ with $C^{-1} = \text{diag}(\alpha)$:

\[
\mathcal{E}(C, \sigma^*) = \int P(Y | X, W, C, \sigma^*) P(W | C) dW
\]

Taking derivatives with respect to hyperparameters yields sparse solutions through Automatic Relevance Determination (ARD).

The updates are:
\[
\alpha_i^{\text{new}} = \frac{\gamma_i}{w_i^2}, \quad (w_i^{\text{new}})^T = \frac{\gamma_i(X_i - w_i^T X)}{N - \sum (1 - E_w[w_i]) \alpha_i}
\]

These sparse solutions effectively perform variable selection during learning.
\end{remark}

\subsection{Bayesian Prediction}

\begin{remark}[Prediction Averaging]
Rather than using point estimates of parameters, Bayesian prediction averages over the posterior:

For unsupervised learning:
\[
P(x | D, m) = \int P(x | \theta, m) P(\theta | D, m) d\theta
\]

For supervised learning:
\[
P(y | x, D, m) = \int P(y | x, \theta, m) P(\theta | D, m) d\theta
\]

The integral naturally favours simpler models through an Occam-factor, since complex models with low likelihood contribute less to the average.
\end{remark}

\subsection{Bayesian Nonparametrics}

\begin{remark}[Nonparametric Models]
Bayesian nonparametric models have no explicit parameters after marginalization. For instance, a Gaussian process defines a prior over functions directly.
\end{remark}

\subsection{Integrating Out Linear Parameters}

\begin{remark}[Linear Regression with Integrated Weights]
If $w \sim N(0, \tau^2 I)$ and $\psi_i \sim N(W^T X_i, \sigma^2)$, integrating $w$ yields:
\[
y | X, D \sim N(\bar{w}^T X, X \Sigma_w X^T + \sigma^2)
\]

If we integrate $w$ with a kernel representation:
\[
\begin{bmatrix} y \\ y_* \end{bmatrix} \bigg| X, x \sim \mathcal{N}\left(\begin{bmatrix} 0 \\ 0 \end{bmatrix}, \begin{bmatrix} K_{xx} & K_{xx_*} \\ K_{x_*x} & K_{x_*x_*} \end{bmatrix}\right)
\]

The conditional predictive distribution is:
\[
y | Y, X, x \sim \mathcal{N}(K_{xx_*} K_{xx}^{-1} Y, K_{x_*x_*} - K_{x_*x} K_{xx}^{-1} K_{xx_*})
\]

This explains the regression prediction through kernel properties.
\end{remark}

\subsection{Posterior Integration as Joint Gaussian}

\begin{remark}[Gaussian Process Conditioning]
Posterior integration with a kernel is equivalent to conditioning a joint Gaussian with that kernel (parameter-free approach).
\end{remark}

\subsection{Covariance Kernel Learning}

\begin{definition}[Covariance Kernel K]
For a covariance kernel $K$, the conditional distribution is:
\[
y | X, K \sim \mathcal{N}(0, K_{xx})
\]

Optimization proceeds via gradient ascent in $\log P(X | X, K)$.
\end{definition}

\newpage

\section{Gaussian Processes and Latent Variable Models}

\subsection{Gaussian Processes}

\begin{definition}[Gaussian Process (GP)]
A Gaussian process is a stochastic process on $\mathbb{R}$ indexed by $x \in \mathcal{X}$, where any finite subset of which have Gaussian distributions. A GP is fully specified by a kernel function $k(\cdot, \cdot)$ and mean function $m(\cdot)$:
\[
f(x) \sim \text{GP}(m(\cdot), k(\cdot, \cdot))
\]

For a set of points $\{x_1, \ldots, x_n\}$, the function values form a vector:
\[
\mathbf{f} = [f(x_1), f(x_2), \ldots, f(x_n)]^T \sim \mathcal{N}(m(x), K_{xx})
\]

where $k(\mathbf{x}_i, \mathbf{x}_j)$ are the covariance kernel evaluations.
\end{definition}

\subsection{Brownian Motion as a Gaussian Process}

\begin{definition}[Brownian Motion]
Brownian motion is a Gaussian process with kernel function:
\[
k(s, t) = \sigma^2 \min(s, t)
\]

where $\sigma^2$ is the variance coefficient per unit time (diffusion coefficient) of the process.
\end{definition}

\subsection{Gaussian Process Regression}

\begin{remark}[GPR Model and Inference]
Standard regression using a GP prior:

\begin{itemize}
    \item \textbf{GP prior}: $f(\cdot) \sim \text{GP}(0, k(\cdot, \cdot))$ implies $f(\mathbf{x}) \sim \mathcal{N}(0, K_{xx})$
    
    \item \textbf{Observation model}: $y_i | x_i, f \sim \mathcal{N}(f(x_i), \sigma^2)$, so $Y \sim \mathcal{N}(0, K_{xx})$ with $\tilde{K}_{xx} = K_{xx} + \sigma^2 I$
    
    \item \textbf{Posterior on latent $f$}: Also a GP:
    \[
    f(\cdot) | X, Y \sim \text{GP}\left(K_{x*}(\tilde{K}_{xx})^{-1} Y^T, k(\cdot, \cdot) - K_{x*}(\tilde{K}_{xx})^{-1} K_{*x}\right)
    \]
    
    \item \textbf{Prediction with observation noise}:
    \[
    y | X, Y, x \sim \mathcal{N}\left(K_{xx'} \tilde{K}_{xx}^{-1} Y, K_{x*x*} - K_{x*x} \tilde{K}_{xx}^{-1} K_{xx*} + \sigma^2\right)
    \]
\end{itemize}
\end{remark}

\subsection{Stationary Gaussian Processes}

\begin{definition}[Translation-Invariant Kernels]
A Gaussian process with translation-invariant covariance kernel is stationary: if $f(t) \sim \text{GP}(0, k)$, then $f(\cdot - x) \sim \text{GP}(0, k)$ for any fixed $x$. This means the covariance depends only on the relative distance between points.
\end{definition}

\subsection{Radial Kernels}

\begin{definition}[Radial and Radially Symmetric Kernels]
A covariance kernel is radial or radially symmetric if:
\[
k(x, x') = h(\|x - x'\|)
\]

A GP with a radial covariance kernel exhibits stationarity with respect to translations, rotations, and reflections of the input space, capturing symmetry in the underlying data-generating process.
\end{definition}

\subsection{Independent Components Analysis (ICA)}

\begin{definition}[ICA Model]
ICA decomposes observed data as a linear mixture of independent non-Gaussian components:
\[
X = \Lambda Z + \varepsilon
\]

where $X \in \mathbb{R}^d$, $Z \in \mathbb{R}^k$ contains non-Gaussian independent components, and $\varepsilon$ is typically zero or small noise. The model is identifiable even when $k > d$ (overcomplete) or $k \leq d$ (undercomplete) provided components are sufficiently non-Gaussian, making it tractable where PCA fails.
\end{definition}

\subsection{Noiseless ICA and Infomax}

\begin{remark}[Noiseless (Infomax) ICA]
In the noiseless case, $x = \Lambda z$, we use the unmixing matrix $W = \Lambda^{-1}$ to recover:
\[
z = W x
\]

The likelihood incorporates the change-of-variables Jacobian:
\[
P(x | W) = \prod_k p_k(W_k x) |W|, \quad P(z) = \prod_k p_k(z_k)
\]

The log-likelihood becomes:
\[
\log P(x) = \log |W| + \sum_k \log p_k(W_k x)
\]

Standard gradient descent update: $\Delta W \propto \nabla_W \log P(x) = (W^T)^{-1} + \nabla_W \sum \log p_k(W_k x)$

A better approach uses natural/covariant gradient which is more efficient for this manifold:
\[
\Delta W \propto \nabla_W \log P(x)(W^T W) = W + g(z) z^T W
\]

Note: EM cannot be used since there is no randomness (posterior is a Dirac delta function), and $A$ is screened from $S$ in the M-step, making it circular.
\end{remark}

\subsection{Mutual Information}

\begin{definition}[Mutual Information]
The mutual information between random variables $X$ and $Y$ measures their statistical dependence:
\[
I(X; Y) = \text{KL}[p_{xy} \| p_x p_y] = H(X) - H(X|Y) = H(Y) - H(Y|X)
\]

where $H$ denotes entropy, measuring the average reduction in uncertainty about one variable from observing the other.
\end{definition}

\subsection{Conditional Entropy}

\begin{definition}[Conditional Entropy]
The conditional entropy of $Y$ given $X$ is:
\[
H(Y|X) = -\sum_x \sum_y p(x,y) \log \frac{p(x,y)}{p(x)}
\]

This measures the average remaining uncertainty in $Y$ after observing $X$.
\end{definition}

\subsection{Feedforward Models and Infomax Learning}

\begin{remark}[Infomax Principle for Learning]
For a feedforward model $\tilde{z} = f_i(z_i)$ where $f_i$ is a squashing function (e.g., sigmoid: $f_i(-\infty) = 0$, $f_i(+\infty) = 1$), the optimal learning objective is:
\[
\underset{W}{\text{argmax}} \, I(X; Z_i) = \underset{W}{\text{argmax}} \left[H(Z) - H(Z|X)\right]
\]

Because $f(W; x)$ is deterministic (given $x$, $z$ is fully determined), we have $H(Z|X) = 0$.

The optimization thus becomes:
\begin{itemize}
    \item \textbf{Maximize} $H(Z)$: Make the output distribution as uniform as possible over $[0, 1]$
    \item \textbf{Constraint}: The sources must become statistically independent as possible
\end{itemize}

This simultaneously maximizes information transmission through the network and encourages statistical independence in the latent representation.
\end{remark}

\newpage

\section{Advanced Models: Boltzmann Machines}

\subsection{Kurtosis and Non-Gaussianity}

\begin{remark}[Kurtosis as Non-Gaussianity Measure]
Kurtosis measures the non-Gaussianity of a source:
\[
\text{Kurtosis} = \frac{E[(x-\mu)^4]}{E[(x-\mu)^2]^2} - 3
\]

A key insight from the Central Limit Theorem: linear mixtures of independent non-Gaussian sources tend toward a more Gaussian distribution, driving kurtosis toward 0. This property makes kurtosis useful for identifying directions that maximize non-Gaussianity in ICA.
\end{remark}

\subsection{Extended Kalman Filter (EKF)}

\begin{definition}[Extended Kalman Filter]
For nonlinear state-space models with nonlinear dynamics and observation functions, the Extended Kalman Filter approximates the posterior by linearizing around the current estimate:

\begin{align*}
Z_{t+1} &= f(Z_t^*, u_t) + \frac{\partial f}{\partial z}\bigg|_{\hat{z}_t} (Z_t - \hat{Z}_t^*) + W_t \\
X_t &= g(\hat{Z}_t^*, u_t) + \frac{\partial g}{\partial z}\bigg|_{\hat{z}_t} (Z_t - \hat{Z}_t^*) + V_t
\end{align*}

The linearization enables approximate inference in nonlinear systems, with Kalman filter-style alternation between prediction and update steps.
\end{definition}

\subsection{Boltzmann Machines}

\begin{definition}[Boltzmann Machine]
An undirected graphical model over binary variables $S_i \in \{0, 1\}$ where some variables are hidden and some are observed:
\[
P(S | W, b) = \frac{1}{Z} \exp\left[\sum_{i < j} W_{ij} S_i S_j - \sum_i b_i S_i\right]
\]

where $W_{ij} = W_{ji}^T$. The energy function is:
\[
H(S) = -\sum_i W_{ij} S_i S_j + \sum_i b_i S_i \quad \text{such that} \quad P(S) = \frac{1}{Z} e^{-H(S)}
\]

Inference requires computing moments under both clamped (observed) and unclamped (full model) distributions:
\[
\langle S^m P(S | S^v, W, b) \rangle \quad \text{and} \quad \langle S^H C H^T P(S | S^v, W, b) \rangle
\]

Learning via EM maximizes:
\[
\log P(C^N H^v | W, b) = \sum \left[W_{ij}^T S_i S_j - \sum b_i S_i - \log Z\right]
\]

The M-step for weight updates separates into clamped expectations $\langle S_i S_j \rangle_c$ (from observed data) and unclamped expectations $\langle S_i S_j \rangle_u$ (from the full model).
\end{definition}

\subsection{Restricted Boltzmann Machines (RBM)}

\begin{definition}[RBM Architecture]
A bipartite Boltzmann machine where $W_{ij} = 0$ for any two visible nodes or any two hidden nodes. This factorization enables efficient inference:

\[
P(S^v | S^H) = \prod_{i \in v} \text{Bernoulli}\left(\sigma\left(\sum_{j \in H} W_{ij} S_j - b_i\right)\right)
\]

\[
P(S^H | S^v) = \prod_{j \in H} \text{Bernoulli}\left(\sigma\left(\sum_{i \in v} W_{ij} S_i - c_j\right)\right)
\]

where $\sigma$ is the sigmoid function. Unclamped samples can be generated efficiently using block Gibbs sampling by alternately sampling visible and hidden layers.
\end{definition}

\subsection{Latent Dirichlet Allocation (LDA)}

\begin{definition}[LDA Model for Topic Modeling]
LDA is a generative model for document collections with latent topics:

\begin{itemize}
    \item \textbf{Objects}: $W$ possible words, $D$ documents, $K$ topics
    
    \item \textbf{Topic distribution per document}: $\theta_d | \alpha \sim \text{Dirichlet}(\alpha, \ldots, \alpha)$ (length $K$)
    
    \item \textbf{Word distribution per topic}: $\phi_k | \beta \sim \text{Dirichlet}(\beta, \ldots, \beta)$ (length $W$)
    
    \item \textbf{Topic assignment}: $z_{id} | \theta_d \sim \text{Discrete}(\theta_d)$ for the $i$-th word in document $d$
    
    \item \textbf{Word identity}: $x_{id} | z_{id}, \phi_{z_{id}} \sim \text{Discrete}(\phi_{z_{id}})$ given the assigned topic
\end{itemize}

The probability of word $w$ in document $d$ integrates over all topics:
\[
P(x_w = w | d) = \sum_{k=1}^K \theta_{dk} \phi_{kw}
\]

This model captures both document-topic and topic-word distributions, enabling unsupervised discovery of latent topics.
\end{definition}

\newpage
\vspace*{2cm}
\begin{center}
    \section*{\Large Part II: Approximate Inference Methods}
\end{center}
\vspace*{1cm}
\nobreak

When exact inference becomes intractable (as with loopy graphs or continuous variables), we resort to approximate methods. The following sections develop the three main paradigms: variational inference, expectation propagation, and loopy belief propagation. All fundamentally rely on minimizing or bounding divergences between the true and approximate posteriors.

\newpage

\section{Variational Inference Methods}

\subsection{Variational Approximation Setup}

\begin{definition}[Variational Approximation in EM]
In variational approximation, if $P(Z | X, \theta^m) \in \mathcal{Q}$ where $\mathcal{Q}$ is a tractable class of approximations, then $\theta^m$ is a fixed point of the variational algorithm.
\end{definition}

\subsection{Variational Lower Bound}

\begin{definition}[Free Energy and ELBO]
The variational free energy is:
\[
\mathcal{F}(q, \theta) = \langle \log p(x, z | \theta) \rangle_{q(z)} + H(q) = \log P(X | \theta) - \text{KL}[q \| p(z | x, \theta)]
\]

The E-step and M-step become:
\begin{itemize}
    \item \textbf{E-step}: $q^*(z) = \arg\max_q \mathcal{F}(q, \theta)$ (optimize q given current $\theta$)
    \item \textbf{M-step}: $\theta^* = \arg\max_\theta \mathcal{F}(q, \theta)$ (optimize $\theta$ given current q)
\end{itemize}
\end{definition}

\subsection{Factored Variational Inference}

\begin{remark}[Factored Variational E-step]
The most common factorization assumes $\mathcal{Q} = \{q : q(z) = \prod_i q_i(z_i)\}$ where the $z_i$ form disjoint sets.

In the $i$-th step of the E-phase, we maximize $\mathcal{F}(q, \theta)$ with respect to $q_i(z_i)$ only, keeping other $q_j$ and parameters fixed:

\[
q_i^*(z_i) = \arg\max_q \mathcal{F}(\ldots)
\]

This iterates until convergence. The update rule is:
\[
q_i(z_i) \propto \exp \left\langle \log p(x, z | \theta) \right\rangle_{\prod_{j \neq i} q_j}
\]

This coordinate-ascent approach ensures the free energy monotonically increases at each update.
\end{remark}

\subsection{Mean-Field Approximations}

\begin{definition}[Mean-Field Factorization]
Mean-field approximations assume complete factorization: $q(z) = \prod_i q_i(z_i)$ where each $z_i$ is updated independently. Each variable receives only its marginal posterior expectation from its neighbours.
\end{definition}

\subsection{Mean-Field for Boltzmann Machines}

\begin{remark}[Mean-Field Updates in BMs]
For a Boltzmann machine with $P(z, s) = \frac{1}{Z} \exp\left(\sum w_{ij} s_i s_j + \sum b_i s_i\right)$, the mean-field update for binary variables is:

\[
\log P(x, z | \theta)_{\text{mean}} = \frac{1}{2} w_{ij} \langle s_i \mathbf{s} \langle s_j \rangle + \frac{1}{2} b_i \langle s_i \rangle
\]

For $S_i \in \{+1, -1\}$ coding, we have $S_i = \sum W_{ij} \langle s_j \rangle$ (treating expectations as fixed parameters).
\end{remark}

\subsection{Mean-Field Approximation for Factorial HMMs}

\begin{remark}[Factored HMM Variational Updates]
For systems with multiple independent chains:
\[
q(S_t^{(m)}) = \prod_m q_t^{(m)}(S_t^{(m)})
\]

The update for chain $m$ is:
\[
q_t^{(m)} \propto \exp \left[ \left\langle \log P(S_t^{(m)}, X_{1:T}) \right\rangle_{\prod_{m' \neq m} q_t^{(m')}} \right]
\]

This decomposes inference into parallel updates across independent components, dramatically reducing computational complexity compared to the full factorial HMM posterior.
\end{remark}

\subsection{Structured Variational Inference}

\begin{remark}[Structured Mean-Field for Complex Models]
Rather than fully factored approximations, structured variational inference uses more complex factorizations that capture important dependencies:

\[
\psi(S_t^{(m)}) = \prod_m \psi^{(m)}(S_t^{(m)})
\]

Each component $\psi^{(m)}$ can capture local dependencies within that chain while remaining independent between chains:
\[
\psi^{(m)}(S_t^{(m)}) \propto \exp \left[ \sum_\ell \log P(S_t^{(m)} | S_{\ell}^{(m)}) + \sum P(X_t | S_t^{(m)}) \right]
\]

This balances computational tractability with improved approximation quality.
\end{remark}

\subsection{Variational Inference on DAGs}

\begin{remark}[Message Passing for DAG Variational Inference]
For directed acyclic graphs $P(X, z) = \prod P(V_k | \text{pa}(V_k))$ with factored assumption $Q(z) = \prod q_i(z_i)$ for disjoint sets $\{z_i\}$:

\[
\log q_i^*(z_i) = \left\langle \sum_k \log P(V_k | \text{pa}(V_k)) \right\rangle_{q_m(z)} + \text{const}
\]

This separates into:
\[
= \sum \langle \log P(z_i | \text{pa}(z_i)) \rangle_{q_m(z)} + \sum \langle \log P(Q_j | \text{pa}(j)) \rangle
\]

where the first sum covers parent messages and the second covers children and their parents. Each node receives messages from its Markov boundary, enabling efficient message-passing inference algorithms.
\end{remark}

\subsection{Variational Bayes}

\begin{definition}[Variational Bayes Algorithm]
Variational Bayes extends mean-field to jointly optimize over both latent variables and parameters:

\[
\log P(x | \mathcal{M}) = \log \int\int P(x, z | \theta, \mathcal{M}) P(\theta | \mathcal{M}) dz d\theta
\]

Using the factorization $Q(z, \theta) = Q_z(z) Q_\theta(\theta)$, the evidence lower bound (ELBO) is:
\[
\mathcal{F}(Q_z, Q_\theta) = \int\int \log \frac{P(x, z, \theta | \mathcal{M})}{Q_z(z) Q_\theta(\theta)} dz d\theta
\]

The variational algorithm alternates:
\begin{itemize}
    \item \textbf{E-like step}: $Q_z^*(z) \propto \exp \langle \log P(z | x, \theta) \rangle_{Q_\theta(\theta)}$
    
    \item \textbf{M-like step}: $Q_\theta^*(\theta) \propto \exp \langle \log P(x, z | \theta) \rangle_{Q_z(z)} P(\theta)$
\end{itemize}

The gap between the true log-evidence and the ELBO is:
\[
\log P(x) - \mathcal{F}(Q_z, Q_\theta) = \text{KL}[Q_z(z) Q_\theta(\theta) \| P(z, \theta | x)]
\]
\end{definition}

\newpage

\section{Complete Exponential Family VB and Advanced Methods}

\subsection{Complete Exponential Family VB}

\begin{remark}[Complete Exponential Family Variational Bayes]
When $Q_0(z)$ is also computed in the complete exponential family:
\[
Q_0(z) \propto P(z) \exp \langle \log P(Z, X | \Theta) \rangle_{Q_z}
\]

This can be written as:
\[
= \int Q(z) \ell(z)^Y \exp \left( q(\theta) / \Gamma \right) + q(\theta)^T \exp \left( \sum_{ij} T(Z_i, X_j) \right)_Q
\]

\[
\propto \int Q(z, \Sigma) q(\theta)^Y \exp \left( q(\theta) / \Gamma \right)
\]

where we only need to track $\bar{b}$ and $\bar{\Sigma}$ with:
\[
\bar{\Sigma} = \bigcup \sum, \quad \bar{\Sigma} = Z + \sum_{ij} \langle T(Z_i, X_j) \rangle_{Q_z}
\]

Similarly:
\[
S_{22}(\hat{Z}) = \int Q_{02}(\hat{Z}) \cdot Q_{02}(\hat{Z}) \times 2 \exp \langle \log P(Z_i; X_i | b) \rangle \times \int q_u(z) \exp \left( \sum_{ij} b_{ij} S T(Z_i, X_j) \right)
\]

\[
= P(\bar{Z}_i | X_i, \bar{\theta}(\theta))
\]

This changes the natural parameter from $q(\theta)$ to $\langle b(\theta) \rangle_{\theta_Q}$.
\end{remark}

\subsection{EM for MAP vs. Variational Bayesian EM}

\begin{remark}[Comparison of EM Objectives]
Two key variants of EM differ in their optimization goals and assumptions:

\textbf{EM Optimization:}
\begin{itemize}
    \item Goal: Maximize $P(D | Z, \theta)$ with respect to $\theta$
    \item E-step: Compute $Q_0(Z) \leftarrow P(Z | X, \theta)$
    \item M-step: $\theta \leftarrow \arg\max \int \log P(Z | Z, \theta; \theta) Q_0(Z) d Z$
\end{itemize}

\textbf{Variational Bayesian EM:}
\begin{itemize}
    \item Goal: Maximize bound on $P(Z | m)$ with respect to $Q_0$
    \item VB-E-step: Compute $Q_0(Z) \leftarrow p(Z | X, \bar{\theta})$
    \item VB-M-step: $\theta_0(\theta) \leftarrow \exp \int \log p(Z, x, \theta) Q_0(Z) dZ$
\end{itemize}

Note: Reduction to EM occurs if $Q_0(\theta) = \delta(\theta - \theta^*)$, in which case $\langle \phi(\theta) \rangle_{\theta(\theta)} = \phi(\theta^*)$.
\end{remark}

\subsection{Variational Free Energy with Hyperparameters}

\begin{remark}[Free Energy Optimization with Hyper-parameters]
The variational free energy can be extended to include hyperparameter optimization:
\[
F(\theta_0, \theta_0, \eta) = \int\int Q_0(Z) Q_0(\theta) \log P(X, Z, \theta | D) dZ d\theta \leq P(Z | \eta)
\]

The hyper-M step maximizes the current bound with respect to $\eta$:
\[
\eta) = \arg\max \int\int Q_0(Z) Q_0(\theta) \log P(Z, \theta, \theta | D) dZ d\theta
\]
\end{remark}

\subsection{ARD with Variational Bayesian Methods}

\begin{remark}[Automatic Relevance Determination]
ARD is used with variational Bayesian methods for determining latent dimensionality.

Example: For $\gamma_n \sim (\Lambda Z, \psi)$ with prior $M'(0,1)$ and column-wise prior $\Lambda_{ii} \sim N(0, \alpha_i^{-1})$:

\[
F(Q_0, Q_n, \psi, \alpha) = \langle \log P(Y | \Lambda, X) \rangle + \log P(\Lambda | \alpha) + \log P(\Psi) \rangle_{Q_\alpha Q_\nu} + \text{const}
\]

Optimization step:
\[
\alpha \leftarrow \arg\max \langle \log P(\Lambda | \alpha) \rangle_{Q_\alpha}
\]

The hyperparameters $\alpha$ are optimized to cause some divergences as ARD, enabling automatic determination of relevant dimensions.
\end{remark}

\subsection{Variational Sparse GP Approximations}

\begin{definition}[Sparse Variational GP Methods]
Variational sparse GP approximations address the computational cost of Gaussian process inference, which scales as $O(N^3)$ for $N$ data points.

The key idea is to select a smaller set of possible function measurements $U$ at input locations $Z$ such that:
\[
P(y | Z, U, x') \approx P(y | x, x', y)
\]

\textbf{GP Notation:}
\begin{itemize}
    \item $F$: latent GP values for $X$
    \item $U$: latent GP values for $Z$
    \item GP prior: $\begin{bmatrix} F \\ U \end{bmatrix} \sim \mathcal{N}\left(0, \begin{bmatrix} K_{xx} & K_{xz} \\ K_{zx} & K_{zz} \end{bmatrix}\right)$
\end{itemize}

\textbf{Marginal Likelihood:}
\[
P(y | X) = \int P(y | F) P(F | U, Z) P(U | Z) dF dU
\]

This factorization enables inducing variable methods for efficient inference in GP models.
\end{definition}

\newpage

\section{Sigma-Point Filtering and Expectation Propagation}

\subsection{Sigma-Point Filter}

\begin{remark}[Sigma-Point Filter for Time Series]
For the state-space model $P(z_{t+1} | x_{t+1})$ with mean $\hat{Z}_t$ and variance matrix $V_t^T$:

\begin{enumerate}
    \item Evaluate the function at multiple sigma points: $f(\hat{Z}_t)$, $f(\hat{Z}_t + \sqrt{V})$
    \item For $K$ dimensions, $\lambda V$ satisfies $V_t^T V = \lambda V$
    \item Fit a Gaussian to the $2K+1$ sigma points to estimate:
    \[
    \hat{Z}_t = f(z_{t+1}) + \Sigma
    \]
    
    \item The fitted Gaussian has parameters:
    \[
    \mathcal{N}\left( \frac{1}{2K+1} \sum f(\delta r), \frac{1}{2K+1} \sum f(\delta r) f(\delta r)^T - \mu \mu^T + \Sigma \right)
    \]
\end{enumerate}
\end{remark}

\subsection{Alternative KL Divergence Minimization}

\begin{remark}[The Other KL Direction]
In some cases, we minimize a different KL divergence:
\[
\arg\min_{Q} \text{KL}[P(z|x) \| \prod T_z(z|x)^T] = \arg\min -\int p(z|x) \log \prod T_z(z|x) dz
\]

This simplifies as:
\[
= \arg\min - \sum \int P(z|x) \log \psi_j(z_j|x) dz = \arg\min - \int P(z|x) \log \psi_z(z|x) dz
\]

which gives us a cross-entropy form (with $-H(\Sigma) + L(p\pi)$), resulting in:
\[
= p(x|z)
\]
\end{remark}

\subsection{Expectation Propagation (EP)}

\begin{definition}[Expectation Propagation Framework]
In EP, we assume the posterior factorizes:
\[
P(Z|X) = \frac{1}{Z_n} \prod P(Z_i | X_i) \propto \prod_i f_i(z_i)
\]

where $z_i$ are not necessarily disjoint. We approximate with sites $\phi(z) \propto \prod_i f_i(z_i)$.

EP minimizes:
\[
\text{KL}[D\mathcal{P} \| \prod_i] \Rightarrow \min_{\phi} \text{KL}\left[\prod f_i(Z_i) \prod f_i(Z_j) \bigg\| \prod f_i^{\text{new}}(Z_i) \prod f_i(Z_j)\right]
\]

Note: If $z_i$ are truly disjoint, we don't need $\prod f_{jz_i}$. EP reduces to a local KL problem.
\end{definition}

\subsection{EP Algorithm}

\begin{remark}[Expectation Propagation Algorithm]
Given sites $f_1(z_1), \ldots, f_W(z_W)$, initialize $\psi_i = \delta_{\text{proj KL}[z_i \| f(z_i)]}$ and $f^{\text{new}}_j = [\text{for } t > 0]$.

\textbf{Algorithm:}
\begin{itemize}
    \item \textbf{Repeat} for $i = 1, \ldots, N$ do:
    \begin{itemize}
        \item \textbf{Delete}: $\psi_i(z) = \frac{\phi(z)}{\hat{\ell}_i(z)} = \prod f_j(z)$
        \item \textbf{Project}: $f_i^{\text{new}}(z_i) = \arg\min_F \text{KL}[f_i(z_i) \psi_i(z_i) \| f_i^{\text{new}}(z_i)]$
        \item \textbf{Include}: $\psi(z) \leftarrow f_i^{\text{new}}(z_i) \psi_i(z)$
    \end{itemize}
    \item \textbf{Until convergence}
\end{itemize}
\end{remark}

\subsection{Message Form in EP}

\begin{definition}[Message Representation in Expectation Propagation]
In EP, messages can be expressed as:
\[
\psi_i(z_i) := \prod_{j \in \text{trees}} M_{j \to i}(z_j \land z_i)
\]

For trees as conditional on $z_i$, the rest form sub-trees that are disjoint and thus independent, so we compute their product. For loopy graphs, this is only an approximation.
\end{definition}

\subsection{Moment Matching in Exponential Families}

\begin{remark}[Projection onto Exponential Families]
If both $\psi_i(z)$ and $\tilde{\psi}$ are in the same exponential family, with $\phi(x) = \frac{1}{Z(\theta)} e^{T(x) \theta}$:

\[
\arg\min_\theta \text{KL}[\phi(x) \| \tilde{\phi}(x)] = \arg\min_\theta \text{KL}\left[\phi(x) \bigg\| \frac{1}{Z(\theta)} e^{T(x)\theta}\right]
\]

\[
= \arg\min_\theta -\int p(x) \log \frac{1}{Z(\theta)} e^{T(x)\theta} dx = \arg\min_\theta - \int p(x) T(x) \cdot \phi dt + \log Z(\theta)
\]

Taking the derivative with respect to $\theta$:
\[
\frac{\partial}{\partial \theta} = -\int dx \, p(x) T(x) + \frac{1}{Z(\theta)} \frac{\partial}{\partial \theta} \int e^{T(x)\theta} dx = -\langle T(x) \rangle_p + \langle T(x) \rangle_q
\]

This is known as \textbf{moment matching}: the parameters are chosen so that the marginal expectations of sufficient statistics match.
\end{remark}

\newpage

\section{Gaussian Process Approximations and Advanced EP Methods}

\subsection{Normal Gaussian Process}

\begin{definition}[Gaussian Process with Non-Gaussian Likelihood]
Consider the Gaussian process $Y(g) X \sim \mathcal{N}(f(x), \sigma^2 I)$:

\[
P(y_1, \ldots, y_n, y_i, \ldots, y_n) = \mathcal{N}(y_1, \ldots, y_n | 0, K) \prod_i \mathcal{N}(y_i | g_i, \sigma_i^2)
\]

where the prior is the GP ($f, f_o(\sigma)$) and likelihood is $\prod_i f_i(y_j)$.

If $f(g_j)$ is not Gaussian, approximate by a Gaussian:
\[
\tilde{y}(g) \propto f_o(\sigma) \prod \hat{f}_i(y_j) \quad \hat{f} \sim \mathcal{N}(0, K t \text{ diag}[C[\hat{f}_j^2]])
\]

The approximate posterior covariance is:
\[
\begin{bmatrix} \mu_{ni} \\ g_j \end{bmatrix} \sim \mathcal{N} \left( \begin{bmatrix} 0 \\ 0 \end{bmatrix}, \begin{bmatrix} \Sigma_{nn0} & \Sigma_{nri} \\ \Sigma_{in} & K_{ir} \end{bmatrix} \right)
\]

Note: $\Sigma_{n,i}$ can be written as $K_{n,i}$ since they are equal. The covariance is $\text{cov}(\mu_{n,j}, y_j) = \text{cov}(g_n + \Sigma_{ni}, y_j) = K_{n,j}$.

Therefore:
\[
\psi_{ij}(g_j) = P(g_j | \tilde{\mu}_{ij}) = \mathcal{N}(\tilde{\Sigma}_{n|i} \tilde{\Sigma}_{nn}^{-1} \tilde{\mu}_{n}, K_i - \Sigma_{i,i}(K_i + \text{diag}(\tilde{\Sigma})_{i,i})^{-1} \Sigma_{i,i})
\]

For EP:
\[
\arg\min_\phi \text{KL}[f_i(y_j) \tilde{\psi}_{ij}(g_j) \| f(g_j) \tilde{\psi}_{ir}(g_j)] \xrightarrow{\text{moment matching}} \hat{f}_i'(g_j) = \frac{\mathcal{N}(\hat{h}(g_j))}{\psi_{ij}(g_j)}
\]

with variance:
\[
\text{var} = \int f(g_j) \tilde{\psi}_{ij}(g_j) g_j^2 dg - (\bar{\mu}^{\text{new}})^2
\]

\textbf{Prediction:}
\[
P(y'|x', \mathcal{D}) = \int P(y'|g') \mathcal{N}(g' | K_{y'x}(K_{xx} + \text{diag}(\hat{\sigma}_i^2))^{-1} K_{xw} - K_{xx}(K_{x'x'} + \text{diag}(\hat{\sigma}_i^2))^{-1} K_{xw}) dg'
\]
\end{definition}

\subsection{Exponential Family Factorization}

\begin{remark}[EP Update for Exponential Family Sites]
If $f_i(z_i) \propto e^{T(z)\theta_i - \psi(\theta_i)}$ (sites in exponential family), then:

\[
q(z) \propto \prod \hat{f}_i^* e^{T(z) \sum \theta_i - Z(\psi(\theta_i))}
\]

Normalizing, we get:
\[
U(z) = e^{T(z) \sum \theta_i - \psi(\sum \theta_i)}
\]
\end{remark}

\subsection{Unnormalized Log-Sum Approximation}

\begin{definition}[Unnormalized Log-Sum Sites]
For approximate sites where $\sum \sum = \sum \theta_i e^{T \Theta \theta}$ with $\theta = \sum \theta_{ij}$ as the natural parameters of $p(z)$:

\[
\tilde{\theta}_i = \sum_j \hat{C}_j \quad \text{for} \quad \psi_{ij}(z), \quad \tilde{L}(\theta) = \log \int e^{T \phi \phi} dz \quad \text{(log normalizer)}
\]
\end{definition}

\subsection{Unnormalized KL Divergence}

\begin{remark}[Unnormalized KL with Moment Matching]
The unnormalized KL divergence is:
\[
\text{KL}[\int p \| q] = \int p(x) \log \frac{p(x)}{q(x)} dx + \int (q(x) - p(x)) dx
\]

This matches the zero moment as well.
\end{remark}

\subsection{Update Rule for Unnormalized Sites}

\begin{definition}[Computation of Site Updates]
For the update:
\[
\int \sum \theta e^{T \Theta \theta} \prod \hat{C}_j e^{T(\Theta)\theta_j} = \int f_i(z_i) \prod \hat{C}_j e^{T(\Theta)\theta_j} dz \Rightarrow \hat{C}_j = e^{\hat{B}_i(\theta_n) - E(\theta)}
\]

We have:
\[
\int f_i(z_i) \prod e^{T(\Theta)\theta_j} dz = e^{\hat{B}_i(\theta_\pi)}
\]

where $\hat{B}_i(\theta_\pi)$ is the log normalizer of the fitted exponential family $\hat{f}_i(z_i) \propto f_i(z_i) e^{T(\Theta)\theta_i}$.
\end{definition}

\subsection{Approximate Log-Likelihood}

\begin{remark}[Log-Likelihood Approximation in EP]
The approximate log-likelihood is:
\[
\log \int \prod f_i(z_i) \approx \log \int \prod \hat{f}_i(z_i) = \log \left( \prod \hat{C}_j \int e^{T(\Theta)\theta} \right) = E(\theta) + \sum \log \hat{C}_j =: \tilde{L}
\]

Since $p(z|x) = \frac{1}{Z} \prod f_i(z_i) = \frac{p(x)}{p(x)} \prod f_i(z_i)$, we have:
\[
p(x) = \int \prod f_i(z_i) dz
\]

Therefore, $\tilde{L}$ is an approximate log-likelihood.
\end{remark}

\subsection{Gradient with Respect to Hyperparameters}

\begin{definition}[Hyperparameter Gradient Computation]
Let $f_i$ depend on model (hyper) parameters $\theta'$:

\[
\nabla_{\theta'} \tilde{L} = \nabla_{\theta'} E(\theta) + \sum \nabla_{\theta'} \log \hat{C}_j \quad \mu = \nabla_\theta E(\theta)
\]

This expands to:
\[
= \mu \nabla_{\theta'} \theta + \sum \nabla_{\theta'} \log \hat{C}_j^* = \mu \nabla_{\theta'} \theta + \sum (\nabla_{\theta'} \hat{B}_i(\theta_\pi) - \mu \nabla_{\theta'} \theta)
\]

\[
= \mu \nabla_{\theta'} \theta + \sum \left( \mu_i \nabla_{\theta'} \theta_n + \langle \nabla_{\theta'} \log f_i(z_i) \rangle_{\hat{\phi}_i} - \mu \nabla_{\theta'} \theta \right)
\]

For the $\hat{B}_i$ gradient:
\[
\nabla_{\theta'} \hat{B}_i(\theta_{\pi i}) = \nabla_{\theta'} \log \{ f_i(z_i) e^{T(\Theta)\theta} \} = \frac{\partial \theta_{\pi i}}{\partial} \hat{B}_i(\theta_{\pi i}) \cdot \nabla_{\theta'} \theta_{\pi i} + e^{\hat{B}_i(\theta_{\pi i})} / \int (\nabla_{\theta'} f_i(z_i)) e^{T(\Theta)\theta} dz
\]

This simplifies to:
\[
= \mu \cdot \nabla_{\theta'} \theta_{\pi i} + \langle \nabla_{\theta'} \log f_i(z_i) \rangle_{\hat{\phi}_i}
\]

where $\langle T(z) \rangle_{\hat{p}_i} = \hat{\mu}_i$ in EP.
\end{definition}

\subsection{Alpha Divergence}

\begin{definition}[Alpha Divergence and Robust Inference]
The alpha divergence is defined as:

\[
D_\alpha[P \| Q] = \frac{1}{\alpha(1-\alpha)} \int (\alpha P(x) + (1-\alpha)Q(x) - P(x)^\alpha Q(x)^{1-\alpha}) dx
\]

Taking limits:
\[
\lim_{\alpha \to 0} D_\alpha(P \| Q) = \text{KL}[Q \| P] \quad \text{and} \quad \lim_{\alpha \to 1} D_\alpha(P \| Q) = \text{KL}[P \| Q]
\]

For robust EP with alpha divergence:
\[
\text{From} = \arg\max KL[f_i(z_i)^\alpha \hat{f}_i(z_i)^\alpha \tilde{\psi}_{ir}(z_j) \| f(z_j) \tilde{\psi}_n(z_j)]
\]

Note: For $\alpha < 1$, we have stable results with small updates from $\hat{f}_i(z_i)$.
\end{definition}

\newpage

\section{Loopy Inference and Loopy BP Methods}

\subsection{Dealing with Loops}

\begin{remark}[Damping and Grouping for Loopy Graphs]
When the graphical model contains loops, two strategies can improve BP accuracy:

\begin{itemize}
    \item \textbf{Damping}: Update messages with damping factor:
    \[
    M_{ij}'(x_j) = (1-\alpha) M_{ij}'(x_j) + \alpha \left( \sum f(x_i, x_j) f(x_j) \prod_{k \to j} M_{k \to j}(x_j) M_{j \to i}(x_j) \right)
    \]
    
    \item \textbf{Grouping variables into groups}: Aggregate variables into groups to improve accuracy
\end{itemize}
\end{remark}

\subsection{Loopy BP as Message-Based Expectation Propagation}

\begin{definition}[Loopy BP as Message-based EP]
Loopy BP can be interpreted as message-based expectation propagation where we approximate pairwise factors by products of messages:

\[
\psi_{ij}(x_i, x_j) = M_{j \to i}(x_j) M_{i \to j}(x_j)
\]

The full joint is approximated as:
\[
P(x) \approx \frac{1}{Z} \prod_{\text{nodes}} f_i(x_i) \prod_{(x_i, x_j)} \frac{f_{ij}(x_i, x_j)}{1/Z'} = \prod_{\text{nodes}} \left( f_i(x_i) \prod_j \frac{M_{j \to i}(x_i)}{M_{j \to i}(x_j)} \right) = \prod_{\text{nodes}} b_i(x_i)
\]

\textbf{EP Interpretation:}
\begin{itemize}
    \item \textbf{Deletion}: Remove the pairwise factor:
    \[
    \psi_{ij}(x_i) = f_i(x_i) f_i(x_j) \prod_k M_{k \to j}(x_j) \prod M_{k \to j}(x_j) \prod \psi_k(x_j) \prod M_{k \to j}(x_k)
    \]
    
    \item \textbf{Projection}: Find the best message approximation:
    \[
    \psi_{ij}(x_i, x_j) = \int M_{j \to j}^{\text{new}}, M_{i \to i}^{\text{new}} = \arg\min \text{KL}[\psi_{ij}(x_i, x_j) \psi_{ij}(x_i, x_j) \| M_{i \to j}(x_j) M_{j \to i}(x_j) \hat{c}_{ij}(x_j)]
    \]
\end{itemize}

The updated beliefs for variable nodes are:
\[
b_i(x_i) \propto M_{j \to i}^{\text{new}}(x_i) \hat{c}_{ij}(x_j) \propto \sum \left( f_{ij}(x_i, x_j) f_i(x_j) \prod_{k \to j} M_{k \to j}(x_j) \right) f(x_j) \prod M_{k \to j'}(x_i)
\]

Thus:
\[
M_{ij}^{\text{new}}(x_i) = \sum \left( f_{ij}(x_i, x_j) f_i(x_j) \prod_{k \to j} M_{k \to j}(x_j) \right)
\]

\textbf{Important Note:} In the tree case, $b_i(x_i) \propto p(x_i)$ is the true marginal. Message-based EP is re-expressed in terms of messages $M_{i \to j}(x_j)$ instead of sites $f_i(z_i)$, so we don't need to approximate both site and cavity. The cavity is approximated as:
\[
\psi_i(z_i) \approx \prod M_{j \to i}(z_j \land z_i) \quad \text{(a product of exponential form)}
\]
\end{definition}

\subsection{Loopy BP as Tree-Based Reparameterization}

\begin{definition}[Tree-based Reparameterization of Loopy BP]
Loopy BP can also be understood as iteratively reparameterizing the original factorization:

\[
P(z) = \frac{1}{Z} \prod_{\text{nodes}} f_i(x_i) \prod_{ij} f_{ij}(x_i, x_j) = P(x_n, y) \prod_{ij} p(x_i | x_j) = \prod_i p(x_i) \prod_{ij} \frac{p(x_i | x_j \epsilon_m(x_j))}{p(x_i) p(x_j)}
\]

where we transition from an undirected tree to a DAG representation using parent marginals where $\sum_{x_i} p(x_i, x_j) = p(x_j)$.

\textbf{Reparameterization Process:}

From $P(z) \propto \prod_{ij} f_{ij}(x_i, x_j)$ to $P(x) = \prod_i p(x_i) \prod_{ij} p(x_i | x_j \epsilon_m(x_j))$, we initialize $f_i^0 = f_j^0 = 1$ and iterate over edges:

\[
P^n(x_i, x_j) = Z_n^0 f_i^n(x_i) f_{ij}^n(x_i, x_j) f_j^n(x_j) \quad \text{(includes } f_i, f_{ij})
\]

Define the updated factors as:
\[
f_{ij}^n(x_i, x_j) = \frac{f_{ij}^n(x_i, x_j)}{M_{j \to i}(x_i)} \quad \text{where we scale down by the message}
\]

The propagated messages ensure consistency across the loopy graph structure. For the marginal computations:
\[
p(x_i) = \prod_i b_i(x_i), \quad \text{where} \quad b_i(x_i) = f_i(x_i) \prod_j M_{j \to i}(x_i)
\]

and:
\[
f_{ij}(x_i, x_j) = \frac{p(x_i, x_j)}{p(x_i)} = \text{(marginalised pairwise)}
\]

\textbf{Pseudo-Marginals in Loopy Graphs:} While the representation ensures loopy messages remain consistent within the spanning tree, the pseudo-marginals are generally not globally consistent.

For pseudo-marginals to be correct:
\[
\sum_{x_j} b(x_i, x_j) = b(x_i), \quad \text{but} \quad \sum_{x_j} \left( b(x_i) \prod_k \frac{b(x_i, x_j)}{b(x_i) b(x_j)} \right) \neq b(x_i)
\]

\textbf{Result:} These pseudo-marginals provide valid beliefs for any of the spanning trees in the loopy graph.
\end{definition}

\newpage
\vspace*{2cm}
\begin{center}
    \section*{\Large Part III: Theoretical Foundations --- Convex Optimization and Duality}
\end{center}
\vspace*{1cm}
\nobreak

The approximate inference methods of Part II can be understood through the lens of optimization theory. The Bethe free energy provides a principled objective for loopy belief propagation, while convex duality reveals the fundamental structure underlying all variational inference procedures.

\newpage

\section{Bethe Free Energy and Convex Duality}

\subsection{Bethe Free Energy on Loopy Consistent Beliefs}

\begin{definition}[Loopy-Consistent Beliefs and Bethe FE]
Loopy BP operates on beliefs that satisfy marginal consistency on trees (not globally):

\[
\sum_{x_i} b_i(x_i) = 1, \quad \forall i \quad \text{and} \quad \sum_{x_j} b_{ij}(x_i, x_j) = b_i(x_i), \quad \forall i, j \in u(i), x_i
\]

The variational free energy is:
\[
F(I) = \sum \log P(x_i) - \Sigma + H[\Sigma]
\]

The \textbf{negative Bethe free energy} is defined as:
\[
F_{\text{bethe}}(b) = E_{\text{bethe}}(b) + H_{\text{bethe}}(b)
\]

where:
\[
P(Y) \propto \prod_i f_i(x_i) \prod_{ij} f_{ij}(x_i, x_j) \propto \prod_i b_i(x_i) \prod_{ij} \frac{b_{ij}(x_i, x_j)}{b_i(x_i) b_j(x_j)}
\]

\textbf{Energy term:}
\[
E_{\text{bethe}}(b) = \sum_i \sum_{x_i} b_i(x) \log f_i(x_i) + \sum_{ij} \sum_{x_i, x_j} b_{ij}(x_i, x_j) \log f_{ij}(x_i, x_j) \quad \text{(using } b_i \text{ as if true marginal)}
\]

\textbf{Entropy term:}
\[
H_{\text{bethe}}(b) = \sum_i H[b_i] - \sum_{ij} \text{KL}[b_{ij} \| b_i b_j] = -\sum_{ij} \sum_{x_i, x_j} b_{ij}(x_i, x_j) \log \frac{b_{ij}(x_i, x_j)}{b_i(x_i) b_j(x_j)}
\]

On a tree, $F_{\text{bethe}} = -F$.
\end{definition}

\subsection{Lagrangian Optimization of Bethe FE}

\begin{remark}[Bethe Free Energy Lagrangian]
The Bethe free energy Lagrangian is:
\[
L = E_{\text{bethe}}(b) + Z_{\text{bethe}}(b) + \left[ \sum a_{ij}(x_i)(-1) + \sum \left( a_{ij}(x_i) \left( \sum_j b_{ij}(x_i, x_j) - b_i(x_i) \right) \right) \right.
\]

\[
\left. + \sum z_{ij}(x_j) \left( \sum_j b_{ij}(x_i, x_j) - b_i(x_j) \right) \right]
\]

Taking derivatives:
\[
\frac{\partial L}{\partial b_i(x)} = \log f_i(x_i) - \log b_i(x_i) + \sum \sum \frac{b_{ij}(x_i, x_j)}{b_i(x_i)} + \sum_j - \sum \sum_{j \in \varepsilon(x)} a_{ij}(x) = 0
\]

\[
\Rightarrow b_i(x_i) \propto f_i(x_i) \prod_{j \in \varepsilon(i)} e^{-z_{ij}(x_i)}
\]

For pairwise beliefs:
\[
\frac{\partial L}{\partial b_{ij}(x, y)} = \log f_{ij}(x_i, x_j) - \log b_j(x_i, x_j) + \log b_i(y) b_j(y_j) + z_{ij}(x_i) + z_{ij}(y_j) = 0
\]

\[
\Rightarrow b_{ij}(x_i, x_j) \propto f_{ij}(x_i, x_j) b_i(x_i) b_j(y_j) e^{z_{ij}(x_i)} e^{z_{ij}(x_j)}
\]

Comparing to BP: $M_{j \to i}(x_i) = e^{-z_{ij}(x_i)}$.

Solving for $z_{ij}(x_i)$ by enforcing the consistency constraint $\sum_j b_{ij}(x_i, x_j) = b_i(x_i)$, we have:

\[
e^{-z_{ij}(x_j)} \propto \sum_{x_j} f_{ij}(x_i, x_j) b_j(x_j) e^{z_{ij}(x_j)} = \sum_{x_j} f_{ij}(x_i, x_j) f_j(x_j) (\varepsilon(x_i)) e^{-z_{ij}(x_j)}
\]

This gives the BP message passing rule.
\end{remark}

\subsection{Fixed Points of Loopy BP}

\begin{definition}[Bethe FE and Loopy BP Fixed Points]
Fixed points of loopy BP are exactly the stationary points of the Bethe tree energy.
\end{definition}

\subsection{Bethe FE as Second-Order Approximation}

\begin{remark}[Inclusion-Exclusion Approximation]
The Bethe free energy is an approximation based on inclusion-exclusion until second order:

\[
H_{\text{bethe}}(b) = \sum_i H[b_i] - \sum_{ij} \text{KL}[b_{ij} \| b_i \otimes b_j]
\]

The Bethe free energy accounts for interactions between different sites, while variational free energy assumes independence.
\end{remark}

\subsection{Log Partition Function}

\begin{remark}[Log Partition Function Properties]
The log partition function $\Sigma(\theta)$ in exponential family $P(x|\theta) = \exp(\theta^T S(X) - \Sigma(\theta))$ has key properties:

\[
\frac{\partial}{\partial \theta} \Sigma(\theta) = \mu, \quad \frac{\partial^2}{\partial \theta^2} \Sigma(\theta) = \text{Var}(S(X))
\]

Thus $\Sigma(\theta)$ is convex in $\theta$. The Kullback-Leibler divergence becomes:

\[
\frac{\partial \theta}{\partial} = \mathbb{E}_\theta[\theta^T S(X) - \Sigma(\theta)] = \theta^T \mu - \Sigma(\theta)
\]
\end{remark}

\subsection{Negative Entropy and Duality}

\begin{definition}[Negative Entropy as Conjugate]
The negative entropy is defined as:

\[
\psi(\mu) = \mathbb{E}_\theta[\log P(X|\theta, \psi)] = \mathbb{E}_\theta[e^{\theta^T S(X) - \Sigma(\theta)}] = \theta^T \mu - \Sigma(\theta)
\]

Thus:
\[
\theta^T \mu = \Sigma(\theta) + \psi(\mu)
\]

The negative entropy is dual to the log-partition function.
\end{definition}

\subsection{Legendre-Fenchel Conjugate}

\begin{remark}[Fenchel Conjugate and Young's Inequality]
The Legendre-Fenchel conjugate is defined as:

\[
f^*(y) = \sup_{y, x} \langle y, x \rangle - f(x) \|_{y,x}
\]

This is always convex, and Young's inequality states:
\[
\langle y, x \rangle \leq f(x) + f^*(y), \quad \forall x, y
\]

The Fenchel-Young inequality shows equality holds when:
\[
f(x) + f^*(y) = x^T y \Longleftrightarrow y \in \partial f(x)
\]

Therefore:
\[
\mu = \nabla \Sigma(\theta), \quad \theta = \nabla \psi(\mu)
\]
\end{remark}

\newpage

\section{Convex Duality and Discrete MRF Inference}

\subsection{KL Divergence and Duality}

\begin{definition}[KL Divergence Bounds and Dual Forms]
Consider the KL divergence:
\[
\text{KL}[q \| p'] = \text{KL}[P(x|\theta_0) P(x|\theta_0)] = \mathbb{E}_\theta[-\log P(x|\theta_0) + \log P(x|\theta)] = -\theta_0^T \mu + \Sigma(\theta_0) + \psi(\mu) \geq 0
\]

Therefore:
\[
\psi(\mu) = \sup_\theta [\theta^T \mu - \Sigma(\theta)]
\]

and continuous functions are reciprocally dual. Thus:
\[
\Sigma(\theta) = \sup_\mu [\theta^T \mu' - \psi(\mu')]
\]
\end{definition}

\subsection{Coupled Exponential Family Representations}

\begin{remark}[Factorization in Exponential Family]
For $p(x,z) = \exp(\theta^T S(x,z) - \Sigma_{xz}(\theta))$ that is also exponential with changed sufficient statistics, the marginal likelihood is:

\[
L(\theta) = p(x|\theta) = \sum_z e^{\theta^T S(x,z)} - \Sigma_{xz}(\theta) = \sum_z e^{\theta S_z(z^x)} e^{-\Sigma_{xz}(\theta)} = \exp(\Sigma_z(\theta) - \Sigma_{xz}(\theta))
\]

The optimization problem becomes:
\[
\mu_z^* = \arg\max_\mu [\theta^T \mu_z - \psi(\mu_z)]
\]

with the feasibility constraint:
\[
\mu = \sum_x \nu(x) S_m(x), \quad \sum_x \nu(x) = 1
\]
\end{remark}

\subsection{Discrete Markov Random Fields}

\begin{definition}[MRF in Exponential Family Form]
A discrete MRF can be written as:

\[
P(x) = \frac{1}{Z} \prod_i f_i(x_i) \prod_{ij} f_{ij}(x_i, x_j) = \exp \left( \sum_i \sum_k \theta_i(k) S_i(x_i=k) + \sum_{ij} \sum_{k,l} \theta_{ij}(k,l) J(x_i=k)(x_j=l) - \Sigma(\theta) \right)
\]

\textbf{Natural parameters:} $\theta = [\theta_i(k), \theta_{ij}(k,l), \forall i,j,k,l]$

\textbf{Mean of sufficient statistics:} $\mu = [P(x_i=k), P(x_i=k, x_j=l), \forall i,j,k,l]$

For MRF with tree structure:
\[
\psi(\mu_T) = \mathbb{E}_\theta \left[ \log \prod_i P(x_i) \prod_{ij} \frac{P(x_i, x_j)}{P(x_i) P(x_j)} \right] = -\sum_i H(x_i) + \sum_{ij} I(x_i, x_j)
\]

where the first sum is entropy and the second sum is mutual information.
\end{definition}

\subsection{Non-Convexity of Bethe Entropy}

\begin{remark}[Non-Convexity Issue]
The Bethe entropy is not convex:
\[
\psi_{\text{bethe}}(\mu_z) = -\sum_i H(x_i) + \sum_{ij} I(x_i, x_j)
\]

This presents challenges for optimization.
\end{remark}

\subsection{Spanning Tree Decomposition}

\begin{definition}[Convex Bound via Spanning Trees]
Consider a set of spanning trees $T$ for the MRF, each with its own parameters $\theta_T, \mu_T$. By padding entries corresponding to off-tree edges with zero, we assume $\theta_T$ has the same dimension as $\theta$.

Suppose we have a distribution $\phi$ on spanning trees such that $\mathbb{E}_\theta[L_{\sigma T}] = \theta$. Then:
\[
\Sigma(\theta) = \tilde{\Sigma}(\mathbb{E}_\theta[L_{\sigma T}]) \leq \mathbb{E}_\theta[\tilde{\Sigma}(\theta_T)]
\]

We can obtain a bound to get a good approximation to $\tilde{\Sigma}$:
\[
\tilde{\Sigma}(\theta) \leq \mathbb{E}_\theta[\Sigma(\theta_T)] = \Sigma_\theta(\theta)
\]

\textbf{Lagrangian formulation:}
\[
\mathcal{L} = \mathbb{E}_\theta[\tilde{\Sigma}(\theta_T)] - \lambda^T (\mathbb{E}_\theta[\theta_T] - \theta) = \Sigma_\theta(\theta)
\]

Taking derivatives:
\[
\frac{\partial \mathcal{L}}{\partial \theta_T} = 0 \Rightarrow \mu_T = \Pi_T(\lambda)
\]

where $\Pi_T(\lambda)$ selects multipliers corresponding to elements of $\theta$ that are non-zero in tree $T$.

\textbf{Dual form:}
\[
\tilde{\Sigma}_\theta(\theta) = \sup_\lambda \left[ \lambda^T \theta + \mathbb{E}_\theta[\Sigma - \psi(\Pi_T(\lambda))] \right]
\]

\[
= \sup_\lambda \left( \lambda^T \theta + \mathbb{E}_\theta[-\psi(\Pi_T(\lambda))] \right)
\]

\[
= \sup_\lambda \left( \lambda^T \theta + \sum H_\lambda(x_i) - \sum_{ij} I_\lambda(x_i, x_j) \right)
\]

where singleton is always selected in $T$. This is the \textbf{convexified version of Bethe free energy}.
\end{definition}

\subsection{REINFORCE Trick}

\begin{remark}[Gradient Estimation via REINFORCE]
The REINFORCE trick provides a way to estimate gradients:

\[
\nabla_\theta F(p, \theta) = \nabla_\theta \int q(z;p) (\log P(z, x|\theta) - \log \mathcal{U}(z;p)) dz
\]

\[
= \int \nabla_\theta q(z;p) (\log P(z, x|\theta) - \log \mathcal{U}(z;p)) dz + \int q(z;p) \nabla_\theta (\log P(z, x|\theta) - \log \mathcal{U}(z;p)) dz
\]

\[
= \int \nabla_\theta \log q(z;p) (\log P(x, z|\theta) - \log \mathcal{U}(z;p)) \mathcal{U}(z;p) dz
\]

\[
= \langle \nabla_\theta \log \mathcal{U}(z;p) (\log P(x, z|\theta) - \log \mathcal{U}(z;p)) \rangle_{\mathcal{U}(z;p)}
\]
\end{remark}

\subsection{Recognition Models and Amortised Inference}

\begin{definition}[Amortised Variational Inference]
Instead of learning an optimal variational distribution $\psi(z)$ that depends on $x$ globally, we learn a \textbf{recognition model}:

\[
\psi(x_i; \theta) \quad \text{where } \theta \text{ is parametrized to depend on } x_i
\]

The optimal variational distribution $\psi^*(z)$ becomes:
\[
B \text{ recognition model, and this approach is called amortised inference}
\]

We learn $f(x_i, \theta)$ instead of $\psi$.
\end{definition}

\newpage
\vspace*{2cm}
\begin{center}
    \section*{\Large Part V: Modern Approaches --- Neural Networks and Generative Models}
\end{center}
\vspace*{1cm}
\nobreak

Recent advances in machine learning combine neural networks with the approximate inference methods developed throughout this course. The following sections cover parameterized variational inference with deep generative models and contemporary approaches like diffusion models.

\newpage

\section{Deep Learning and Parameterized Variational Inference}

\subsection{Wake-Sleep Learning Algorithm}

\begin{definition}[Wake-Sleep Learning]
The wake-sleep algorithm alternates between two phases for learning both generative and recognition models:

\textbf{Wake Phase:}
\begin{itemize}
    \item Generative model: $p_\theta(x, z)$ with real data $x$
    \item Infer: $z = f(x; \theta)$ with current $f$
    \item Update generative parameter: $\theta \leftarrow \theta + \alpha \gamma \nabla_\theta \log f(j, \theta)$
\end{itemize}

\textbf{Sleep Phase:}
\begin{itemize}
    \item Recognition model: $q(z|x; f = f(x, \varphi))$, sample $[z, x, z_s]$ from $p(x|z)$
    \item Update to minimize: $\text{KL}[P(z|z) \| q(z, f(x, \varphi))] \Rightarrow \min -\text{avg}_{z,s} \langle \log f(x, \varphi) \rangle$
    \item Gradient update: $\Delta \varphi \propto \sum_s (\nabla_s - f(x, \varphi)) \nabla_\varphi f(x, \varphi)$
\end{itemize}

\textbf{Alternation:}
\begin{itemize}
    \item \textbf{Wake:} Freeze $\varphi$, update $\theta$ to explain real data
    \item \textbf{Sleep:} Freeze $\theta$, update $\varphi$ to recover latents on model-generated data
\end{itemize}
\end{definition}

\subsection{Variational Autoencoders (VAE)}

\begin{definition}[Variational Autoencoder Architecture]
A VAE consists of:
\[
P(z) = \mathcal{N}(0, I), \quad P(x|z) = \mathcal{N}(f_{NN}(z; \theta), \sigma^2 I)
\]

where $f_{NN}(z; \theta)$ is Gaussian with a neural network for $\theta$.

\textbf{Posterior approximation (encoder):} $q(z|x; f(x, \varphi))$

\textbf{Variational Objective:}
\[
F(\theta, \varphi) = \left\langle \sum_z \log P(x|z) + \log P(z) - \log q(z|x; f(x, \varphi)) \right\rangle = \sum_{\text{data}} \log P(x|z_{\text{data}}) - \text{KL}[q(z|x) \| P(z)]
\]

\[
= - \sum_{\text{data}} \left\langle \frac{\|x - g(z)\|^2}{2\sigma^2} \right\rangle_z - \text{KL}[q(z|x) \| P(z)]
\]

where the first term is the reconstruction cost and the second is the regulariser.
\end{definition}

\subsection{Reparameterization Trick}

\begin{remark}[Gradient Estimation via Reparameterization]
If $f$ gives marginal distribution with mean $\mu$ and variance $\Sigma_2$ for $z$, we can use:
\[
z_i^S = \mu_i + \sigma_i \varepsilon_i^S, \quad \varepsilon_i^S \sim \mathcal{N}(0, 1)
\]

The objective becomes:
\[
F_{\hat{f}}(\theta, \varphi) = \frac{1}{S} \sum_S \log P(x_i, z_i^S; \theta) - \log q(z_i^S; f(x_i, \varphi))
\]

Gradients are computed as:
\[
\frac{\partial F_{\hat{f}}}{\partial \theta} = \frac{1}{S} \sum_S \nabla_\theta \log P(x_i, z_i^S; \theta)
\]

\[
\frac{\partial F_{\hat{f}}}{\partial \varphi} = \frac{1}{S} \sum_S \frac{\partial}{\partial z_i} [\log P(x_i, z_i^S; \theta) - \log q(z_i^S; f(x_i, \varphi))] \frac{dz_i^S}{d\varphi} + \frac{\partial}{\partial f(x_i)} \log q(z_i^S; f(x_i, \varphi)) \frac{df(x_i)}{d\varphi}
\]

Train $\theta$ and $\varphi$ together, holding $z^S$ fixed.
\end{remark}

\subsection{Importance-Weighted Variational Bounds}

\begin{definition}[Importance-Weighted Free Energy]
The importance-weighted lower bound improves upon the standard ELBO:

\[
\gamma(\psi) = \log P(x) - \log \mathbb{E}_{z \sim p_q(x)} [\text{-} \log P(x)|z] \geq \log \mathbb{E}_z \left[ \frac{P(x,z)}{q(z)} \right] > -F(\psi, \varphi)
\]

With $K$ importance samples:
\[
\mathcal{L}(k) = \log \mathbb{E}_{z_1, \ldots, z_K \sim q} \left[ \frac{1}{K} \sum P(x, z_i) / q(z_i) \right] \geq \mathbb{E}_{z_1, \ldots, z_K \sim q} \left[ \log E/K \sum \frac{P(x, z_i)}{q(z_i)} \right]
\]

As $k \to \infty$, the signal for learning amortised $q$ vanishes as we lose the gradient of $z$:

\[
\nabla_\varphi L_k = \mathbb{E}_{z_1, z_K \sim q} \left[ \frac{1}{K} \sum (w_k / \sum w_i) \nabla_\varphi \log q(z_k) \right], \quad w_k = \frac{P(x, z_k)}{q(z_k)}
\]

where $w_k = \frac{P(x,z_k)}{q(z_k)}$ and $\text{large } k + 1/K \sum w_i \approx \sum K$ contribute to the gradient.

\textbf{Note:} $\nabla_\varphi L_k \propto K \mathbb{E}_\theta[\nabla_\varphi q(z)] = 0$ as $\mathbb{E}[\nabla_\varphi \log q(z)] = 0$.

For small $k$: The signal dominates at some $k$, gradient is strong, and $q(z,z)$ learning focuses on high-likelihood region.
\end{definition}

\subsection{Normalising Flows}

\begin{definition}[Normalising Flows for Flexible Posteriors]
Consider a recognition model $q(z)$ implicitly defined by:
\[
z_0 \sim p_0(\cdot|x) \leftarrow \text{a fixed, tractable on, e.g. } \mathcal{N}(x, I)
\]

Then:
\[
z = f_K(f_{K-1}(\ldots f_1(z_0))), \quad f_k \text{ is smooth, invertible parametrised by } \varphi
\]

The expected log probability transforms as:
\[
\mathbb{E}[F(z)]_{z} = \mathbb{E}[F(f_K(f_{K-1}(\ldots f_1(z_0))))]_{z_0}
\]

and:
\[
\log q(z_0) = \log p_0(f_K^{-1}(\ldots f_1(z_0))) - \sum_k \log \left| \frac{\partial f_k}{\partial z_k} \right|
\]

Since $z_k = f_k(z_{k-1})$:
\[
q(z_{k+1}) = q(f_k^{-1}(z_k)) \left| \frac{\partial z_{k+1}}{\partial z_k} \right|^{-1} = q(f_k^{-1}(z_k)) |\nabla f_k(z_{k-1})|^{-1}
\]

Given a sample $z_0^S \sim q_0(\cdot; x)$:
\[
F(\varphi, \theta) \propto \frac{1}{2} \sum \log p(x_i, f_K(\ldots f_1(z_i^S)); \theta) + H[z_0] + \frac{1}{2} \sum_i \sum_k \log |\nabla f_k(f_{k-1}(\ldots f_1(z_i^S)))|
\]
\end{definition}

\newpage

\section{Diffusion Models and Score-Based Methods}

\subsection{Diffusion Model Overview}

\begin{definition}[Forward and Reverse Diffusion Processes]
A diffusion model combines a forward (noising) process with a learned reverse (denoising) process:

\textbf{Fixed Forward Diffusion Process:}
Starting from data $x_0$, progressively add Gaussian noise:
\[
q(z_t | z_{t-1}) = \mathcal{N}(\sqrt{1-\beta_t} z_{t-1}, \beta_t I)
\]

where $\beta_t$ is a predefined schedule. After $T$ steps:
\[
q(z_T) \approx \mathcal{N}(0, I)
\]

\textbf{Markov Chain Structure:}

The forward process forms a Markov chain: $z_0 \to z_1 \to z_2 \to \cdots \to z_T$, where each step adds noise. The reverse process learns to denoise: $z_T \to z_{T-1} \to \cdots \to z_0$.
\end{definition}

\subsection{Variational Generative Model (Backward Model)}

\begin{definition}[Reverse Process in Variational Framework]
The variational generative model learns to reverse the noising process:

\[
p(z_T) = \mathcal{N}(0, I)
\]

\[
p_\theta(z_{t-1} | z_t) = \mathcal{N}(\mu_\theta(z_t, t), \sigma_t^2 I)
\]

\[
p_\theta(x | z_0) = \mathcal{N}(\mu_\theta(z_0, 0), \sigma_0^2 I)
\]

\textbf{Evidence Lower Bound (ELBO):}

The ELBO decomposes as:
\[
\log p(x) \geq \mathbb{E}_{q(z_{1:T}|x)} \left[ \log p_\theta(x|z_0) + \sum_{t=1}^T \log \frac{p_\theta(z_{t-1}|z_t)}{q(z_t|z_{t-1})} \right]
\]

This is commonly rewritten as a sum of diffusion losses at each timestep $t$. The forward process posterior at step $t$ is:
\[
q(z_t | x) = \mathcal{N}(\sqrt{\bar{\alpha}_t} x, (1-\bar{\alpha}_t) I)
\]

where $\bar{\alpha}_t = \prod_{s=1}^t (1 - \beta_s)$. The variational model learns to approximate this posterior in reverse.
\end{definition}

\subsection{Recognition-Parametrized Model (RPM)}

\begin{definition}[Recognition Models as Variational Posterior Approximations]
Instead of using a fixed variational distribution, we parameterize the approximate posterior (encoder) as a neural network:

\[
q_\phi(z_{t-1} | z_t, x) = \mathcal{N}(\mu_\phi(z_t, t, x), \sigma_\phi^2(z_t, t, x) I)
\]

The parameters $\phi$ of the encoder depend on the observation $x$ and current time step $t$. This creates an amortized inference mechanism where a single encoder network handles all possible observations.

\textbf{Joint Training:}
Both the denoising model $p_\theta$ and the recognition model $q_\phi$ are trained jointly. The recognition model helps guide the reverse process by providing information about the data at each denoising step, improving sample quality and training efficiency compared to unconditional denoising.
\end{definition}

\end{document}
